\section{The Smart Contracts and Their Triggers}\label{sec:contracts}

We have the door, the log, and the declared data a contract may read. What remains
is the contract itself: the thing that, when an event fires, works out what must
change and submits it. This section specifies what a contract is allowed to be,
shows that trades, lifecycle events, and corporate actions are one kind of machine
rather than three, and makes the two-layer correctness statement precise on worked
numbers.

\subsection{One architecture, three triggers}

To the ledger, executing a trade and processing a lifecycle event are the same
operation: a finite list of atomic moves with optional state changes, obeying
conservation, admitted at the one door. There is no separate lifecycle subsystem
inside the ledger, no per-product pipeline, no special data path for corporate
actions. This is not an incidental economy; it is a direct consequence of the
thesis. The ledger is a door, some checks, a log, and views, and nothing that
admits or records a transaction cares where the transaction came from. So
everything that changes the ledger is a transaction at the same door, and
everything that produces a transaction presents the ledger with the same shape: a
proposal to be admitted or refused. The producers' internals differ; the shape they
present does not.

That shape has three trigger classes, and they differ in what they read, not in
how they run:

\begin{itemize}[leftmargin=1.5em]
\item A \textbf{trade} is a contract triggered by execution. The executed trade
becomes one conserving transaction; for a bespoke instrument the executed trade is
itself the unit definition, registered at the moment it is done.
\item A \textbf{lifecycle event} is a contract triggered by \emph{the engine} --- the
external execution service that fires dated obligations --- when an obligation comes
due, an expiry, a coupon, an exercise. It reads the unit's current terms and state
and the observations stamped for it, and it emits the moves that discharge the
event.
\item A \textbf{corporate action} is a contract triggered by a confirmed notice ---
a dividend, a split, a spin-off. It reads the notice's declared parameters and the
unit's schedule, and it submits the boundary's one transaction, which may carry all
three of a new terms version, the entitlement moves, and the status writes that set
the new basis.
\end{itemize}

Each trigger fires a contract that produces exactly one transaction per firing. That
one transaction may carry many moves, a terms version, and several status writes at
once, but it is one atomic commit; a follow-on that a boundary sets in motion --- a
cascade onto a successor unit, a downstream settlement --- is a separate trigger
with its own transaction, not a second output of the first.

An option expiry and a stock split are processed by the same shape of machine:
catch the trigger, read declared data and stamped observations, compute, submit.
The differences between them are entirely differences in what they read. Nothing
in the ledger branches on which of the three a given transaction came from.

The corporate-action trigger is a \emph{confirmed} notice, and the gap between an
announcement and its confirmation is handled deliberately. An announcement enters as
declared data --- the boundary is known, its effective time recorded --- but it arms
no operator: until confirmation, no conversion is declared, and nothing about any
datum adjusts. The unit is marked pending for that boundary. Consumption that would
cross the announced boundary fails closed and quarantines, because there is no
operator to carry a number across it; consumption that stays in the unit's standing
basis --- a fresh print taken before the effective time and read in the frame it was
taken --- is unaffected and proceeds normally. Confirmation is precisely the act
that arms the operator, turning the announced boundary from a known-but-inert marker
into a declared conversion the read seam can apply.

A single door settles quietly what happens when two things land on one unit at one
instant, because there is no such instant in the log: every commit is totally
ordered by the door, so events that share an effective time are simply adjacent, in
a definite order. Two cases are worth stating. When two \emph{boundaries} carry the
same effective time on the same unit, the same-effective-time weld applies and
effective order keys on the triple (effective time, declared precedence, bid); if
the two notices declare no mutual precedence, the second is refused and the unit
enters a pending-transition state --- fail-closed --- rather than letting the bid
tie-break silently choose an order for two adjustments that do not commute. When a
\emph{lifecycle event} and a \emph{boundary} fall together, there is no dedicated
priority scheme, and none is needed: log commit order fixes which prefix each
processor reads, prefix-determinism makes each submission a function of that prefix,
and a settlement booked against a prefix that a boundary later precedes was correct
on the prefix it read: recomputation reproduces it, diffs to zero, and does not flag
it. What the retro-effective path guarantees is that the invalidation is exactly
locatable and repairable, not that it is automatically found. The boundary's window
$[t_{\text{eff}}, t_{\text{notify}})$ is explicit in the log; replaying the corrected
chain recomputes what each settlement in the window should have been; and the repair
is a compensating transaction --- authored outside the ledger, like every repair,
admitted with lineage to both the settlement it corrects and the boundary that
invalidated it. The guarantee is order, determinism, and repairability --- not a
ranking of event types.

A contract, then, is a deterministic, total function from its inputs to one
transaction --- a finite list of moves together with the terms versions and status
writes that ride with them --- with no randomness, no hidden state, and no dependency
on anything not passed to it. Its output alphabet is exactly three things: moves,
which route units between wallets and never create or destroy them; terms versions,
appended never edited; and status writes, each an absolute replacement through its
one legal writer. Conservation holds by construction over the moves, and the terms
versions and status writes pass through the same admission checks as any other
transaction. This is also why writing a contract from a legal confirmation is
transcription rather than engineering. The confirmation already says what happens
on each date and in each state; the contract is its machine-readable form, adding
executability and adding nothing to the information.

\subsection{What the contract layer is required to satisfy}

Because a contract is untrusted, ``deterministic and total'' has to be made
precise enough to check. Four obligations do that, and every implementation of a
contract --- the reference one, a vendor's system, a rewrite --- is held to exactly
them.

\begin{itemize}[leftmargin=1.5em]
\item \textbf{Closed inputs.} A contract reads the committed log and the stamped
observations the log records --- nothing else. No ambient clock, no direct feed, no
configuration off to the side. The timer that wakes a lifecycle contract is
delivery, not data: it says \emph{when}, while everything the event \emph{means}
is in the log. An expiry contract woken an hour late submits the identical
transaction, because nothing it reads changed in that hour. The engine that fires
these timers lives outside the trust boundary, named apart from the door, the log,
and the views, and holding no privileged access of any kind; it is bound only to
fire durably and to deliver at least once. Because its trigger is delivery and never
data, an engine that stalls, fails over, fires early, or fires twice can \emph{delay}
a discharge but cannot corrupt the log: a late delivery submits the same transaction,
a duplicate delivery is absorbed idempotently at the door, and an early or spurious
fire finds its observation or due-state not yet in the log and yields a refusal --- a
rejected proposal, never a corrupt commit.
\item \textbf{Prefix-determinism.} Given a log prefix, the submission is
determined. Two runs, or two independent implementations, over the same prefix
produce equal transactions --- equal move for move, under a \emph{canonical move
order}, the deterministic total order on a transaction's moves fixed by the
transaction itself, against which two recomputations are compared --- so equality is
decided rather than judged.
\item \textbf{Idempotence keys.} Every submission carries a key derived
deterministically from its cause: the obligation's identity for a lifecycle
settlement, the causing boundary paired with the target unit for a corporate-action
follow-on. A retried submission is bit-identical, carries the same transaction id,
and the door commits it once. Crash-and-retry cannot double-book.
\item \textbf{Proposed, never authoritative.} A submission becomes ledger fact
only by passing the same admission checks as every other transaction. No contract
holds write privilege; no submission channel bypasses the door.
\end{itemize}

The reference implementation of a contract is a pure function --- roughly,

\begin{lstlisting}
handle :: LogPrefix -> Trigger -> Either Refusal Transaction
\end{lstlisting}

--- and the four obligations are exactly the properties this signature already
enforces: it takes the log prefix and the trigger and returns a transaction, with
no argument through which a clock, a feed, or a configuration could enter, and no
effect it could have except returning the transaction it proposes. Any other
implementation is conforming precisely insofar as it matches what this function
does on every prefix.

These four obligations make a contract's \emph{output} checkable, but they say
nothing about whether an obligation fires at all. That is a separate guarantee, and
the design carries it as data rather than as hope. Every contractual obligation is a
first-class logged object, registered atomically with the position it binds and
carrying three things: a deadline, a discharge predicate, and a compensation action.
By its deadline it must reach one of three terminal states --- Discharged,
Compensated, or Defaulted --- so a coupon that should fire and does not is not
silence but a \emph{Pending-past-deadline} object, visible to any fold of the log
and therefore a detectable state rather than a missing one. The guarantee is
conditional, and its primary condition is named: it rests on the engine being
available at the deadline to deliver the trigger; Section~\ref{sec:framework} states
the full set of conditions the guarantee stands on. Under an engine outage the
obligation is reached \emph{eventually}, when delivery resumes, not instantaneously;
the log never loses the obligation, it only waits for the fire.
Section~\ref{sec:framework} develops the obligation object and its liveness
guarantee, including which component drives it to a terminal state.

\subsection{Two layers of correctness, on numbers}

We can now discharge the claim from Section~\ref{sec:thesis} that a contract's
output is trusted for nothing and verifiable in everything. The subject is the
contract, not the ledger: the ledger is precisely the trusted component --- it is
trusted to admit only well-formed transactions and to record them faithfully --- and
it is the contracts whose output is trusted for nothing and checked against the log
instead. The claim splits into two layers, and the honest content of the split is
that one layer is guaranteed and the other is checked.

The first layer is structural, and it holds against a defective contract and a
malicious one alike, because it never depended on the contract. Whatever a
contract submits, if the door admits it then it conserves every unit, applies in
full or not at all, writes every fact through its one legal writer, and dangles no
reference. A contract cannot break this by being wrong, because being wrong is not
one of the ways a transaction can be malformed. There is simply no unconserved,
half-applied, or dangling proposal for the door to let through.

The second layer is economic, and this is the layer a contract can violate.
Consider a settlement whose declared terms owe \EUR{100} per note, and a defective
contract that submits a settlement paying \EUR{98}. That transaction is balanced,
conserving, and basis-consistent. The door has no ground to refuse it --- every
structural property holds --- and it is admitted. On a 1{,}000-note holding it has
now booked \EUR{98{,}000} where \EUR{100{,}000} was owed. Nothing at the door
catches this, and nothing at the door is supposed to. What catches it is
recomputation: the settlement is recomputed from the log prefix --- the terms, the
schedule, the registered kinds, the stamped observations are all there --- and
compared with what was recorded, in canonical form. The comparison reports a
difference of \EUR{2} per note, \EUR{2{,}000} on the holding. That nonzero
difference is the defect, located exactly, and it is a defect in the contract, not
in the ledger. It is repaired downstream by a compensating transaction --- authored
outside the ledger, by the corrected producer, and resubmitted through the same door
under a \texttt{corrects} field that names the \EUR{98} settlement it supersedes. The
ledger admits this repair exactly as it admits any transaction and manufactures none
of it; the original \EUR{98} settlement stays in the log, marked corrected, so the
error and its repair both survive.

The corporate-action half of the system runs the same check on the same kind of
numbers. A declared \texttt{Scale 2}, composed with the schedule, takes a
\EUR{100} strike to \EUR{50}. A defective contract that appends \EUR{49} instead
has submitted something structurally admissible --- a well-formed terms version ---
and it is caught the same way, by recomputing the appended strike from the boundary
and the schedule and finding the \EUR{1} gap. One architecture, one verification
mechanism, on both halves: a settlement that should pay \EUR{100} and an adjustment
that should yield \EUR{50} are checked by the same recomputation, and a wrong
\EUR{98} or a wrong \EUR{49} is found the same way.

This is the whole of what Principle~\ref{prin:two-layer} claims, now on numbers.
The ledger \emph{relies on} the structural layer, because it is guaranteed by
construction and no contract can break it. It \emph{expects} the economic layer,
and rather than trusting the expectation it makes it checkable: ``the contract
computed the right number'' is not believed, it is recomputed. A failed
recomputation is always a producer defect and never a ledger defect, because the
ledger's part --- admitting only well-formed transactions and recording them
faithfully --- is exactly the part that was guaranteed.

The component that performs this recomputation deserves a name, because it, and not
the door, is what makes economic correctness checkable. Call it the \emph{witness}.
The witness takes the committed log, replays the declared inputs a transaction
depended on --- the terms, the schedule, the registered kinds, the stamped
observations --- recomputes the transaction in canonical form, and compares. Every
input it reads is log-resident and its computation is deterministic, so any
independent party can run it and reach the identical verdict; who operates it --- the
firm, a counterparty, an auditor, a regulator --- is a deployment choice, not a
matter of trust. The witness is not part of the door. The door admits well-formed
transactions in log order and moves on; the witness reads the committed log after the
fact. Nothing gates admission, a view, or the settlement projection on a prior
witness pass, and when the witness runs is an operational decision the deployment
must schedule. This is why the
bound is eventually-detected-and-repaired rather than verified-before-consumed: the
check is real and anyone can run it, but it runs over what is already committed, not
before it.

\subsection{Provenance is recorded, never acted on}

Every transaction that follows from a boundary carries a tag naming its cause ---
the boundary that triggered it. It is tempting to think the ledger uses this tag to
decide how to process the follow-on. It does not. The tag is provenance only. A
unit's own boundary transaction carries no cause tag, because a boundary is its own
cause: its identity is the transaction's own entry in the log, and a tag pointing a
boundary at itself would be redundant at best and circular at worst. Every follow-on
and cascade, by contrast, carries the causing boundary's identity. But no admission
check, no valuation, no conversion, and no code path anywhere branches on the tag's
value. It is read only by the witness that verifies economic correctness, and by
audit.

The one gate that does look at the tag is a well-formedness check, not a behavioural
branch: a transaction that names a cause is admitted only if that cause is a
boundary already committed strictly earlier in the log. This refuses a transaction
that points at an event not yet recorded, and it structurally rules out a
transaction that names itself. The tag names the cause; it does not carry the
conversion. Audit needs the name, arithmetic needs the operator, and the two are
kept apart on purpose: the tag alone is necessary for provenance and insufficient
for computation, which is exactly why nothing computes from it.

\subsection{The model-agnostic wall}

One boundary has been implicit throughout and is worth making explicit, because it
is what keeps the whole architecture indifferent to product complexity. The ledger
records the outputs of pricing models and never runs a pricing model. Whatever
computation decided that a structured note was worth what it was worth --- however
elaborate, however proprietary --- is not something the ledger reproduces,
validates, or even sees. What the ledger sees is the moves the contract emits and
the observations stamped into the log. The repricing an event causes is new
information, and it enters only as a fresh stamped observation; the ledger never
manufactures it by transforming an old observation into a new one. The gap between
a transported old price and the next fresh print is market profit and loss, and it
is left as such, never absorbed into a boundary.

This is why the four obligations above quantify over the log prefix and the stamped
observations and mention no model. A pricing function is a black box on the far
side of the wall. A declared operator --- the one generic conversion of
Section~\ref{sec:data} --- is a known, closed-menu transformation on the near side.
The wall between them is what lets the ledger process a bespoke note and a plain
share with the same machine: it never has to understand either one, only to admit
the conserving transactions their contracts propose and to recompute those
transactions from data that is entirely in the log.
